% !TEX root = ../report.tex

\chapter{Теоретичні основи колоризації зображень}
    \label{chap:review_theory}
        У цьому розділі розглядаються основні теоретичні аспекти колоризації зображень, включно з історією розвитку методів, класичними алгоритмічними підходами та сучасними нейромережевими моделями.
        Особлива увага приділяється архітектурам від згорткових мереж до дифузійних моделей, а також методам оцінки якості отриманих результатів.

        Основна мета цього розділу --- проаналізувати еволюцію методів колоризації від класичних згорткових мереж до сучасних трансформерів,
        виявити їхні обмеження щодо обчислювальної складності та збереження просторових контурів об'єктів,
        і на основі цього обґрунтувати необхідність переходу до гібридних ієрархічних моделей на базі селективних просторів станів.

    \section{Поняття та історія колоризації}

        Колоризація --- це процес додавання кольорової інформації до монохромних зображень.
        З математичної точки зору, цю задачу можна сформулювати як відображення одноканального сірого зображення у триканальне кольорове зображення, наприклад, у просторі RGB, CIELAB або YUV \cite{2025-survey}.
        Оскільки одному значенню яскравості може відповідати безліч кольорових відтінків, задача є невизначеною та має характер <<один до багатьох>>, що вимагає використання апріорних знань про семантику сцени.

        Історично методи колоризації можна поділити на три основні етапи:
        \begin{enumerate}
            \item \textbf{Ручна та напівавтоматична колоризація:} вимагала значних зусиль художників або використання простих евристик на основі сегментації.
            \item \textbf{Методи на основі глибокого навчання (CNN та GAN):} поява великих наборів даних, таких як ImageNet та COCO, дозволила тренувати згорткові мережі для передбачення кольору \cite{2016-deep-learning}.
                    Знаковим став підхід \textit{Colorful Image Colorization} \cite{2016-colorful-image-colorization}, де задачу було переформульовано з класичної регресії на класифікацію.
                    Замість того, щоб мінімізувати усереднену похибку, яка призводить до тьмяних кольорів, модель передбачає розподіл ймовірностей для кожного пікселя по палітрі відтінків.
            \item \textbf{Генеративні та послідовні моделі (Transformers, Diffusion та SSM Models):} сучасний етап,
            що характеризується використанням механізмів глобальної уваги та імовірнісних моделей дифузії для генерації високоякісних текстур,
            а також впровадженням селективних моделей простору станів (SSM) для лінійного моделювання контексту зображень та забезпечення бездоганної семантичної узгодженості \cite{2023-diffusion-models-survey, 2023-mamba}.
        \end{enumerate}

        Розвиток методів глибокого навчання, зокрема поява архітектур U-Net та GAN, дозволив перейти від простого розфарбовування до генерації фотореалістичних зображень \cite{2017-image-to-image-translation}. 
        Останні дослідження все частіше фокусуються на дифузійних моделях, які демонструють state-of-the-art результати у задачах синтезу зображень \cite{2022-palette-image-to-image-diffusion-models}. 
        
        Поряд із цим, активного розвитку набули архітектури на базі механізмів глобальної уваги. Проте висока обчислювальна складність класичних трансформерів та нездатність їхніх лінійних альтернатив зберігати точний глобальний контекст роблять їх неоптимальними для інтерактивних задач.

        Водночас новітні архітектури на базі SSM здатні забезпечити глобальне рецептивне поле з лінійною складністю.
        Хоча вони й почали з'являтися у вузькоспеціалізованих задачах відновлення інфрачервоних \cite{2024-colormamba, 2026-haqagen} та супутникових знімків \cite{2026-mfmamba}, вони залишаються недослідженими в контексті інтерактивної колоризації фотографій.
        Відсутність рішень, які б поєднували лінійну обчислювальну ефективність просторів станів із точним просторовим контролем за допомогою користувацьких підказок, створює простір для наукового пошуку та підкреслює актуальність запропонованого в цій роботі гібридного підходу.
    
    \section{Підходи до колоризації}

        Підходи до колоризації базуються на різних архітектурах нейронних мереж.
        Опустимо гібридні архітектури та розглянемо деякі з базових.

        \subsection{Згорткові нейронні мережі (CNN)}
            Ранні підходи використовували автоенкодери для прямого передбачення каналів кольоровості (наприклад, $a$ та $b$ у просторі CIE Lab).

            Ларсссон та ін. \cite{2016-learning-representations-for-automatic-colorization, 2016-colorful-image-colorization} запропонували використовувати ознаки, вивчені для класифікації об'єктів, для задачі колоризації, демонструючи тісний зв'язок між семантикою об'єкта та його кольором.

            Інший популярний метод, \textit{Deep Koalarization} \cite{2017-deep-koalarization}, використовує Inception-ResNet-v2 для вилучення ознак та подвійний декодер для відновлення кольору.

        \subsection{Генеративно-змагальні мережі (GAN)}
            З метою усунення ефекту недостатньої насиченості кольорів були запропоновані різні GAN.

            \begin{itemize}[label=$\bullet$] 
                \item Модель \textit{Pix2Pix} \cite{2017-image-to-image-translation} використовує умовний GAN (cGAN) для перетворення зображень, що дозволяє генерувати більш чіткі та реалістичні результати.
                \item Метод \textit{ChromaGAN} \cite{2020-chromagan} поєднує геометричну інформацію з семантичним розподілом класів для покращення правдоподібності розфарбовування.
                \item Також існують підходи для непарного навчання, такі як CycleGAN \cite{2017-unpaired-image-to-image-translation}, що дозволяють тренувати моделі без наявності пар <<чорно-біле -- кольорове>>.
            \end{itemize}

        \subsection{Трансформери}
            З появою механізму уваги дослідники почали застосовувати трансформери для колоризації.

            Наприклад, \textit{Colorization Transformer} \cite{2021-colorization-transformer} та \textit{UniColor} \cite{2022-unicolor} використовують архітектуру трансформерів для моделювання довгих залежностей у зображенні, що покращує узгодженість кольорів на великих відстанях.

            Як ще один приклад: \textit{DDColor} \cite{2023-ddcolor} використовує подвійні декодери та запити кольору (color queries) для досягнення фотореалістичності.

        \subsection{Дифузійні моделі}
            Найбільш перспективним напрямком на сьогодні є імовірнісні дифузійні моделі.
            Вони працюють шляхом поступового видалення шуму з випадкового сигналу, керуючись вхідним чорно-білим зображенням.
            \begin{itemize}[label=$\bullet$]
                \item \textbf{Palette} \cite{2022-palette-image-to-image-diffusion-models} --- універсальний фреймворк для перетворення зображень, який перевершує GAN та регресійні моделі без необхідності налаштування під конкретну задачу.
                \item \textbf{Cold Diffusion} \cite{2023-cold-diffusion} пропонує підхід до інверсії довільних перетворень зображень без використання гаусового шуму, що відкриває нові можливості для відновлення зображень.
            \end{itemize}

            % Для відео ключовою проблемою є часова узгодженість (temporal consistency).
            % Сучасні методи, такі як \textit{Stable Video Diffusion} \cite{2023-stable-video-diffusion} та \textit{VanGogh} \cite{2025-vangogh}, адаптують дифузійні моделі для обробки послідовностей кадрів (додають вимір часу, щоб бачити відео цілком).
            % Інші підходи, наприклад \cite{2019-deep-exemplar-based-video-colorization} та \cite{2024-temporally-consistent-video-colorization}, використовують поширення ознак (feature propagation) (механічно переносить кольори з попередніх кадрів, щоб вони не змінювалися) та саморегуляризацію (метод навчання, який штрафує мережу, якщо вона змінює колір одного й того ж об'єкта у сусідніх кадрах) для уникнення мерехтіння.

        \subsection{Моделі простору станів (SSM)}
            Новітнім трендом у глибокому навчанні є впровадження селективних моделей простору станів, зокрема архітектури Mamba \cite{2023-mamba}.
            Вони розроблені як ефективна альтернатива трансформерам, що здатна моделювати наскрізні глобальні контекстні зв'язки зображення,
            але має лінійну обчислювальну складність, що критично для систем реального часу.

    \section{Інтерактивна колоризація та багатоваріантні результати}

        Оскільки колоризація є відносно суб'єктивною задачею, важливу роль відіграють методи, що дозволяють користувачеві впливати на результат.

        \textbf{Колоризація на основі прикладів (Exemplar-based):}
        Методи, такі як \textit{Deep Exemplar-based Colorization} \cite{2018-deep-exemplar-based-colorization}, використовують референсне кольорове зображення для перенесення колірної палітри на цільове чорно-біле.

        \textbf{Колоризація на основі тексту (Language-based):}
        З розвитком мультимодальних моделей, таких як CLIP, з'явилася можливість керувати кольором за допомогою текстових описів.
        \begin{itemize}[label=$\bullet$]
            \item \textbf{L-CAD} \cite{2023-l-cad} використовує дифузійні пріори для колоризації на основі описів будь-якого рівня деталізації.
            \item \textbf{Diffusing Colors} \cite{2023-diffusing-colors} пропонує фреймворк, де текстові підказки керують процесом дифузії, забезпечуючи баланс між автоматизацією та контролем.
        \end{itemize}

        \textbf{Колоризація на основі підказок користувача (User-guided):}
        Методи \textit{Real-Time User-Guided Image Colorization} \cite{2017-real-time-user-guided-image-colorization} та \textit{iColoriT} \cite{2023-icolorit} дозволяють користувачеві ставити кольорові точки або штрихи, які мережа поширює на відповідні семантичні регіони, використовуючи, наприклад, Vision Transformer для кращого розуміння контексту.

        Сучасна тенденція, як показано в \textit{Control Color} \cite{2025-control-color} та \textit{UniColor} \cite{2022-unicolor}, полягає у створенні уніфікованих фреймворків, що підтримують різні типи умов, такі як текст, зразок, точки, штрихи, в одній моделі.

    \section{Метрики якості колоризації}

        Оцінка якості колоризації є складною задачею через відсутність єдиного <<правильного>> розв'язку: одному чорно-білому зображенню може відповідати безліч правдоподібних кольорових варіацій. Тому для комплексного оцінювання розроблених моделей використовують набір різнопланових метрик, які охоплюють як математичну точність, так і перцептивний реалізм та обчислювальну ефективність.

        Використовуються наступні основні групи метрик:
        \begin{itemize}[label=$\bullet$]
            \item \textbf{Метрики цільової похибки ($\mathrm{MSE}$ або $\mathrm{MAE}$):} У задачах інтерактивної колоризації середньоквадратичну або середньоабсолютну похибку доцільно розділяти на дві складові:
                \begin{itemize}[label=$\square$]
                    \item $\mathrm{MSE}_{\mathrm{LAB}} / \mathrm{MAE}_{\mathrm{LAB}}$ --- загальна похибка відновлення кольору в усьому просторі зображення.
                    \item $\mathrm{MSE}_{\mathrm{HINT}}$ --- похибка виключно в тих точках, де користувач надав підказки. Ця метрика є критично важливою для керованих моделей, оскільки значення, близькі до нуля, доводять ідеальну інтеграцію та безумовне дотримання моделлю користувацького вводу без ефекту ігнорування.
                \end{itemize}

            \item \textbf{Піксельні та структурні метрики:} $\mathrm{PSNR}$ (Peak Signal-to-Noise Ratio) та $\mathrm{SSIM}$ (Structural Similarity Index). Вони попіксельно порівнюють згенерований результат з оригінальним кольоровим зображенням.
                Хоча вищі значення вказують на краще відновлення оригіналу, ці метрики часто не корелюють з людським сприйняттям, оскільки <<інший>> правдоподібний колір математично розцінюється як помилка.

            \item \textbf{Перцептивні генеративні метрики:} $\mathrm{LPIPS}$ (Learned Perceptual Image Patch Similarity), $\mathrm{FID}$ (Fréchet Inception Distance) та $\mathrm{KID}$ (Kernel Inception Distance). 
                На відміну від піксельних метрик, вони порівнюють не конкретні пікселі, а високорівневі семантичні ознаки, вилучені за допомогою попередньо навчених мереж, таких як VGG чи Inception.
                Зниження показників $\mathrm{FID}$ та $\mathrm{KID}$ свідчить про те, що розподіл згенерованих кольорів максимально наближений до розподілу реальних фотографій, що підтверджує високу візуальну реалістичність.
                Оскільки метрика $\mathrm{KID}$ обчислюється ітеративно на випадкових підмножинах даних, у роботі вона подається у форматі двох значень:
                \begin{itemize}[label=$\square$]
                    \item $\mathrm{KID}_{\mathrm{mean}}$ --- середнє значення метрики, визначає загальну якість генерації.
                    \item $\mathrm{KID}_{\mathrm{std}}$ --- стандартне відхилення, визначає стабільність та дисперсію результатів моделі на різних вибірках.
                \end{itemize}

            \item \textbf{Обчислювальна та апаратна ефективність:} Для систем, що передбачають інтерактивну роботу користувача, фундаментальне значення мають показники швидкодії.
            До них належать час інференсу ($\mathrm{Time}$, вимірюється в мілісекундах) та обсяг споживаної відеопам'яті ($\mathrm{VRAM}$).
            Сучасні архітектури повинні забезпечувати баланс між високою генеративною якістю та мінімальним споживанням ресурсів графічного процесора.
        \end{itemize}

    % \label{chap:theory}

        % \begin{figure}[h!]
        %     \centering
        %     \includegraphics[width=1\textwidth]{../Preps/StructureScheme.png}
        %     \caption{Структурна схема класифікації методів колоризації}
        %     \label{fig:StructureScheme}
        % \end{figure}

    %     У цьому розділі проводиться більш детальний аналіз сучасних алгоритмів колоризації, що зображені на структурній схемі (рис. \ref{fig:StructureScheme}).
    %     Ми розглянемо еволюцію підходів від класичних оптимізаційних задач до новітніх генеративних моделей.

    %     Основна мета цього розділу --- виявити сильні та слабкі сторони існуючих рішень за такими критеріями: точність відтворення кольорів, обчислювальна складність, можливість керування результатом та часова узгодженість для відеопослідовностей.

    % \section{Алгоритмічні методи колоризації зображень}

    % \section{Методи колоризації на основі нейронних мереж}

    \section{Більш детально про існуючі підходи до колоризації}

        До епохи глибокого навчання колоризація розглядалася переважно як задача оптимізації.
        Класичні алгоритмічні методи можна розділити на дві великі групи: методи на основі втручання користувача (Scribble-based) та методи перенесення кольору зі зразку (Example-based).

        \subsection{Scribble-based}
            Цей підхід базується на припущенні, що сусідні пікселі з подібною інтенсивністю повинні мати подібний колір.
            Користувач наносить на чорно-біле зображення кольорові штрихи, які слугують граничними умовами.
            Задача зводиться до мінімізації цільової функції енергії $J$ \cite{2025-survey, 2004-colorization-using-optimization}:

            \begin{equation*}
                J(U) = \sum_{r} \left( U(r) - \sum_{s \in N(r)} w_{rs} U(s) \right)^2,            
            \end{equation*}
            де $ U(r)$ --- колір у пікселі $r$, $N(r)$ --- окіл пікселя, а $w_{rs}$ --- ваговий коефіцієнт, що залежить від подібності яскравості пікселів $r$ та $s$.

            Перевагою методу є повний контроль користувача, проте він вимагає значних ручних зусиль і часто дає артефакти на межах об'єктів зі слабким контрастом.

        \subsection{Example-based}
            У цьому підході колір переноситься з референсного кольорового зображення на цільове чорно-біле шляхом зіставлення статистик яскравості та текстур.
            
            Ранні роботи \cite{2001-reinhard-color-transfer,2002-welsh-transferring-color} використовували зіставлення гістограм у \textbf{декорельованому} колірному просторі $l\alpha\beta$ (простір Рудермана), що дозволяло маніпулювати колірними каналами незалежно один від одного без виникнення візуальних артефактів.
            
            Хоча цей метод є більш автоматизованим, він критично залежить від вдалого вибору референсного зображення та часто призводить до помилкового забарвлення семантично різних об'єктів, що мають схожу текстуру.

        \subsection{Згорткові нейронні мережі (CNN)}
            Поява CNN дозволила моделювати складні залежності між формою об'єкта та його ймовірним кольором, використовуючи великі набори даних.

            Перші нейромережеві методи намагалися прямо передбачити значення каналів $ab$ у просторі Lab, мінімізуючи середньоквадратичну помилку (MSE):
            \begin{equation*}
                \mathcal{L}_{L2} = \frac{1}{2} \sum_{h,w} || Y_{gt} - \hat{Y} ||_2^2,
            \end{equation*}
            де $Y_{gt}$ --- справжнє зображення, а $\hat{Y}$ --- згенероване.

            Проте, як зазначають Zhang et al. \cite{2016-colorful-image-colorization}, використання MSE призводить до усереднення всіх можливих кольорів, що дає ненасичені результати.

            Для вирішення цієї проблеми задачу було переформульовано як мультимодальну класифікацію, де мережа передбачає розподіл ймовірностей для квантованих відтінків кольору, що значно підвищило яскравість результатів.

        \subsection{Генеративно-змагальні мережі (GAN)}
            Мережі GAN, зокрема архітектура Pix2Pix \cite{2017-image-to-image-translation}, здійснили прорив у чіткості зображень.
            Генератор $G$ намагається створити реалістичне зображення, щоб обдурити дискримінатор $D$, який вчиться відрізняти справжні кольорові фото від згенерованих.
            Функція втрат GAN має вигляд:
            \begin{equation*}
                \mathcal{L}_{cGAN}(G, D) = \mathbb{E}_{x, y}[\log D(x, y)] + \mathbb{E}_{x, z}[\log(1 - D(x, G(x, z)))].
            \end{equation*}
        
            Підходи на кшталт ChromaGAN \cite{2020-chromagan} інтегрують семантичну інформацію у процес генерації, що дозволяє уникнути грубих помилок, наприклад, зеленого неба.

        \subsection{Трансформери та механізми уваги}
            Останні дослідження \cite{2021-colorization-transformer, 2022-unicolor, 2023-ddcolor} показують, що класичні згорткові мережі мають локальне рецептивне поле, через що вони погано справляються з моделюванням глобальних зв'язків. Це часто призводить до семантичних помилок: наприклад, різні частини одного об'єкта (як-от капот і дах автомобіля), перекриті іншими елементами сцени, можуть бути розфарбовані в різні кольори.
            
            Перехід до архітектур на базі трансформерів дозволив вирішити цю проблему. Вони використовують механізм глобальної уваги (Self-Attention), який дозволяє кожному патчу зображення взаємодіяти з усіма іншими патчами одночасно, враховуючи контекст усієї сцени. Це є критично важливим для просторово узгодженого забарвлення великих областей та підтримки єдиної логіки освітлення.
            
            Зокрема, модель Colorization Transformer \cite{2021-colorization-transformer} вперше успішно застосувала цей підхід, розбивши зображення на послідовність патчів і розглядаючи задачу колоризації як задачу прогнозування послідовності токенів (подібно до генерації тексту). Модель генерує кольори авторегресійно, що забезпечує високу якість, проте робить процес інференсу надзвичайно повільним.

            Значним кроком у напрямку інтерактивності стала модель UniColor \cite{2022-unicolor}. Вона запропонувала уніфікований фреймворк, який здатний обробляти мультимодальні умови --- від автоматичної колоризації до генерації на основі текстових описів, референсних зображень або точкових підказок (штрихів). Використання трансформера дозволило ефективно інтегрувати ці різнорідні підказки у єдиний простір репрезентацій.

            Для вирішення проблеми розмиття контурів, яка часто виникає у чистих трансформерів, архітектура DDColor \cite{2023-ddcolor} впровадила систему подвійних декодерів (Dual Decoders). Один декодер відповідає виключно за відновлення просторової роздільної здатності (піксельний рівень), а інший --- за формування кольорових семантичних ознак (глобальний контекст), після чого їхні результати зливаються. Це дозволило досягти високого рівня фотореалістичності.

            Проте, незважаючи на беззаперечні переваги в якості генерації, використання класичного механізму уваги накладає жорсткі апаратні обмеження. Обчислювальна складність Self-Attention зростає квадратично $\mathcal{O}(N^2)$ зі збільшенням роздільної здатності зображення. Це робить такі моделі вкрай вимогливими до обсягу відеопам'яті ($\mathrm{VRAM}$) і ускладнює їх використання для інтерактивної колоризації зображень у високій якості в реальному часі.

            Спроби подолати цю проблему за рахунок так званих «ефективних трансформерів» призводять до нових компромісів. Наприклад, обчислення уваги в локальних вікнах, наприклад, в Swin Transformer \cite{2021-swin-transformer}, або використання математичних лінійних апроксимацій змушують модель втрачати здатність до безкомпромісного та миттєвого глобального моделювання контексту \cite{2022-efficient-transformers-survey, 2021-long-range-arena}.
            У задачі інтерактивної колоризації це стає критичним недоліком, оскільки нейромережа втрачає здатність швидко та безшовно поширювати колір від точкової підказки користувача на віддалені ділянки об'єкта.
            
        \subsection{Дифузійні імовірнісні моделі}
            На відміну від GAN, які генерують результат за один прохід, дифузійні імовірнісні моделі формують зображення ітеративно.
            Їхня математична основа базується на двох марковських процесах:
            прямому, який поступово руйнує дані шляхом додавання гаусового шуму,
            та зворотному, де нейронна мережа вчиться крок за кроком відновлювати оригінальний сигнал,
            керуючись вхідним чорно-білим зображенням як просторовою умовою.

            Цей клас моделей здійснив революцію в синтезі зображень завдяки високій стабільності навчання,
            тобто відсутності проблеми колапсу мод, який характерний для GAN, та неперевершеній якості генерації дрібних текстур.
            Знаковою роботою в цій галузі став універсальний фреймворк Palette \cite{2022-palette-image-to-image-diffusion-models},
            який продемонстрував, що умовна дифузійна модель здатна перевершити спеціалізовані GAN-архітектури у задачі колоризації без додаткового налаштування під специфіку предметної області. 

            Поряд із класичним підходом активно розвиваються альтернативні формулювання генеративного процесу.
            Наприклад, метод Cold Diffusion \cite{2023-cold-diffusion} довів можливість інверсії довільних детермінованих перетворень зображень без використання чистого гаусового шуму,
            що розширює теоретичні межі застосування дифузійних алгоритмів.

            Проте фундаментальним недоліком дифузійних моделей залишається їхня надзвичайно висока обчислювальна вартість на етапі інференсу.
            Оскільки для створення фінального зображення модель повинна виконати десятки або сотні послідовних ітерацій денойзингу, час генерації зростає на порядки порівняно з моделями прямого проходу.
            У контексті керованої інтерактивної колоризації, де фундаментальною вимогою є миттєвий візуальний зворотний зв'язок на кожну нову підказку користувача, використання чистих дифузійних моделей є недоцільним через руйнування користувацького досвіду.
        
        \subsection{Моделі простору станів (SSM) та архітектура Mamba}
            Новітнім класом нейромережевих архітектур, що дозволяє подолати квадратичну складність трансформерів та повільність дифузійних моделей,
            є селективні SSM, зокрема архітектура Mamba \cite{2023-mamba}.
            Вони здатні моделювати наскрізні глобальні контекстні зв'язки з суворою лінійною обчислювальною складністю $\mathcal{O}(N)$,
            що робить їх ідеальними для систем реального часу.

            Математично базова система SSM відображає вхідний одновимірний сигнал $x(t) \in \mathbb{R}^L$ у прихований стан $h(t) \in \mathbb{R}^N$ через лінійні диференціальні рівняння:
            \begin{align*}
                h'(t) &= A \cdot h(t) + B \cdot x(t), \\
                y(t) &= C \cdot h(t) + D \cdot x(t),
            \end{align*}
            де $A, B, C, D$ --- матриці параметрів системи.
            При цьому матриця $A$ характеризує внутрішню динаміку системи та керує еволюцією прихованого стану, тобто відповідає за довготривалу пам'ять,
            матриця $D$ описує прямий зв'язок між входом і виходом, виконуючи роль класичного пропускного з'єднання,
            а матриці $B$ та $C$ визначають коефіцієнти взаємодії прихованого стану із вхідним та вихідним сигналами відповідно.

            Для обробки дискретних послідовностей безперервні параметри трансформують у дискретні ($\overline{A}, \overline{B}$) із використанням кроку дискретизації $\Delta$. Зазвичай для цього застосовують метод фіксації нульового порядку, Zero-Order Hold:
            \begin{align*}
                \overline{A} &= \exp(\Delta A), \\
                \overline{B} &= (\Delta A)^{-1}(\exp(\Delta A) - I) \cdot \Delta B,
            \end{align*}
            що дозволяє переписати систему у дискретному вигляді: $h_k = \overline{A} \cdot h_{k-1} + \overline{B} \cdot x_k$, $y_k = C \cdot h_k + D \cdot x_k$.

            Головна інновація архітектури Mamba полягає у селективності її внутрішніх механізмів:
            параметри $B$, $C$ та крок дискретизації $\Delta$ стають динамічними функціями від вхідного сигналу,
            на відміну від статичних матриць $A$ та $D$.
            Завдяки такій дискретизації параметрів та використанню апаратно-оптимізованого алгоритму паралельного сканування в пам'яті графічного процесора,
            модель динамічно фільтрує нерелевантну інформацію та ефективно обробляє довгі послідовності без обчислювальних обмежень механізму уваги.

            Однак пряме застосування оригінальної одновимірної Mamba до двовимірних зображень призводить до проблеми втрати локального контексту,
            оскільки розгортання матриці пікселів у вектор розриває їхню просторову суміжність по вертикалі.
            Для вирішення цієї проблеми візуальні адаптації на кшталт Vision Mamba \cite{2024-vim} та новітні моделі спектральної трансляції \cite{2024-colormamba} використовують механізми двонаправленого або перехресного сканування.
            Це дозволяє моделі агрегувати ознаки з різних напрямків, повністю відновлюючи просторове розуміння сцени.

            У контексті керованої колоризації використання Mamba відкриває унікальну можливість:
            здатність моделі миттєво та узгоджено поширювати колірну інформацію від точкової підказки користувача на весь об'єкт завдяки глобальному рецептивному полю,
            уникаючи при цьому експоненційного зростання споживання відеопам'яті.
            Саме цей синергетичний баланс обчислювальної ефективності та просторової експресивності робить SSM найбільш перспективним фундаментом для розробки сучасних інтерактивних систем розфарбовування.

    \section{Порівняння різних методів колоризації}
        Узагальнимо характеристики розглянутих методів у порівняльній таблиці \ref{tab:comparison}.
        Для аналізу обрано такі ключові критерії:
        \begin{itemize}[label=$\bullet$]
            \item \textbf{Якість (Quality):} рівень фотореалістичності, природність та насиченість згенерованих кольорових відтінків.
            \item \textbf{Узгодженість (Consistency):} просторова зв'язність кольору в межах об'єктів та відсутність візуальних артефактів.
            \item \textbf{Керованість (Control):} здатність моделі точно інтерпретувати та безумовно інтегрувати підказки користувача.
            \item \textbf{Швидкість (Speed):} обчислювальна складність та час інференсу, що визначають придатність до роботи в реальному часі.
        \end{itemize}

        \begin{table}[ht!]
            \centering
            \caption{Порівняльний аналіз методів колоризації зображень}
            \label{tab:comparison}
            \begin{tblr}{
                width = \textwidth,
                colspec = {X[2,l,t] X[2,l,t] X[3,l,t] X[3,l,t]},
                row{1} = {font=\bfseries, halign=c},
                hlines,
                vlines,
            }
                Метод / Архітектура    & Представники                                                                                                                  & Переваги                                                                                                    & Недоліки                                                                                                                 \\
                Класичні оптимізаційні & Levin et al. \cite{2004-colorization-using-optimization}, Welsh \cite{2002-welsh-transferring-color}                          & Точний локальний контроль користувача, відсутність семантичних <<галюцинацій>>                              & Обчислювально повільні, вимагають значної ручної розмітки, пасують перед складними текстурами                            \\
                CNN                    & Zhang et al. \cite{2016-colorful-image-colorization}, Larsson \cite{2016-learning-representations-for-automatic-colorization} & Висока швидкість інференсу, стабільність процесу навчання                                                   & Обмежене рецептивне поле; схильність до генерації усереднених, тьмяних відтінків при мінімізації $\mathrm{MSE}$          \\
                GAN                    & Pix2Pix \cite{2017-image-to-image-translation}, ChromaGAN \cite{2020-chromagan}                                               & Висока чіткість контурів, насичені та візуально привабливі кольори                                          & Нестабільність змагального навчання, схильність до появи колірних артефактів та плям                                     \\
                Transformers           & ColTran \cite{2021-colorization-transformer}, UniColor \cite{2022-unicolor}, DDColor \cite{2023-ddcolor}                      & Моделювання глобального контексту всієї сцени, гнучка мультимодальна керованість                            & Квадратична обчислювальна складність $\mathcal{O}(N^2)$, критично високі вимоги до обсягу відеопам'яті ($\mathrm{VRAM}$) \\
                Diffusion Models       & Palette \cite{2022-palette-image-to-image-diffusion-models}, Cold Diffusion \cite{2023-cold-diffusion}                        & Найвища якість генерації (SOTA), бездоганний фотореалізм текстур                                            & Екстремально низька швидкість інференсу через сотні послідовних ітерацій денойзингу                                      \\
                Mamba SSM              & Vision Mamba \cite{2024-vim}, ColorMamba \cite{2024-colormamba}                                                               & Суворо лінійна складність $\mathcal{O}(N)$, стабільне споживання $\mathrm{VRAM}$, глобальне рецептивне поле & Чутливість до одновимірного розгортання 2D-даних, ризик втрати або розмиття локальних контурів об'єктів                  \\
            \end{tblr}
        \end{table}

        Проведений аналіз демонструє наявність компромісу між обчислювальною лінійністю та збереженням локальних просторових контурів у сучасних послідовних моделях.
        Чисті трансформери забезпечують високу якість зв'язності, але стають апаратно непридатними для обробки великих зображень у реальному часі, тоді як SSM пропонують бажану лінійну ефективність $\mathcal{O}(N)$,
        але потребують додаткових механізмів локального посилення контурів. 

        Таким чином, гібридні підходи, що поєднують генеративну силу та швидкість селективних лінійних моделей простору станів із точністю локального вилучення ознак за допомогою ієрархічних згорткових декодерів,
        виглядають найбільш обґрунтованими та перспективними для вирішення задачі керованої інтерактивної колоризації.
        Саме розробці та дослідженню такої гібридної архітектури присвячено наступні розділи цієї роботи.

\chapconclude{\ref{chap:review_theory}}

    У цьому розділі було проведено огляд теоретичних основ та сучасних методів колоризації зображень.
    Було показано, що еволюція методів пройшла шлях від простих евристик до складних генеративних моделей, таких як GAN, Transformers, дифузійні фреймворки та селективні моделі простору станів.

    Аналіз літератури свідчить, що в той час як задача безумовної колоризації досягла значних успіхів у реалізмі, існують відкриті питання в області обчислювальної складності моделей при роботі з високою роздільною здатністю, керованості процесом та багатоваріантності для усіх моделей.
    Зокрема, гібридні підходи, мультимодальні підходи та уніфіковані дифузійні й рекурентні структури є найбільш актуальним напрямком досліджень.

    Порівняння нейромережевих підходів показало, що хоча архітектури GAN забезпечують швидкий інференс, вони часто страждають від нестабільності навчання та появи колірних артефактів.
    Дифузійні моделі дозволяють досягти найвищої якості та фотореалізму текстур, проте їхнє практичне застосування жорстко обмежене тривалим часом інференсу через ітеративний характер процесу генерації.
    У свою чергу, підходи на базі трансформерів ефективно моделюють глобальний контекст усієї сцени, однак квадратичне зростання обчислювальних витрат $\mathcal{O}(N^2)$ суттєво ускладнює їх масштабування для зображень високої роздільної здатності.
    
    У наступному розділі буде запропоновано власний підхід до вирішення задачі колоризації зображень, що базується на гібридизації архітектури U-Net та Mamba SSM з використанням спільних ваг.
    Це дозволить ефективно поєднати вилучення низькорівневих просторових ознак, для боротьби з color bleeding, із лінійним моделюванням глобального контексту та суттєвим збереженням пам'яті графічного процесора.

    % Наприкінці кожного розділу ви повинні навести коротенькі підсумки по його 
    % результатах. Зокрема, для оглядового розділу в якості висновків необхідно 
    % зазначити, які задачі у даній тематиці вже були розв'язані, а саме 
    % поставлена вами задача розв'язана не була (або розв'язана погано), тому у 
    % наступних розділах ви її й розв'язуєте.

    % До другого розділу також краще написати малесенький вступ. Зокрема, це 
    % збільшує загальний об'єм роботи та покращує її читабельність.


    % У другому розділі необхідно наводити розв'язання поставленої перед вами 
    % задачі у теоретичному або аналітичному сенсі (хоча, звісно, все залежить 
    % від того, яка саме задача перед вами поставлена).

    % Бажано, щоб кожен пункт завдань, окреслених у вступі, відповідав певному 
    % розділу або підрозділу у дипломній роботі.

    % Посилання до роботи рекомендовано робити за допомогою підсистеми BibTex. 
    % Це максимально зручно, повірте мені. Усі притомні світові журнали надають 
    % посилання на свої статті у вигляді bib-файлів, з яких просто треба 
    % перенести відомості у свій файл. Існують інструменти автоматизованого створення bib-файлів, наприклад, JabRef.
    % Якщо у вас нема автоматизованого механізму, то посилання на джерела у 
    % bib-файлах легко створюються вручну. Головне --- не забувати чотири простих 
    % правила:

    % 1) Імена авторів необхідно подавати у вигляді <<Прізвище, Ім'я>> та 
    % розділяти ключовим словом <<and>>, наприклад:

    % \texttt{author = ''Яковлєв, Сергій Володимирович and Дамблдор, Альбус Персиваль Вульфрік Брайан''}

    % 2) Після усіх полів запису необхідно ставити кому, окрім останнього.

    % 3) Джерела у переліку посилань з'являться тільки після того, як ви зробите 
    % посилання на нього у тексті, наприклад, див. роботи~\cite{ct1, ct2}.

    % Наприкінці розділу знову наводяться коротенькі підсумки.